\documentclass[11pt,twocolumn,a4paper]{article}

% ─── Page layout ─────────────────────────────────────────────────────
\usepackage[margin=2cm, columnsep=0.6cm]{geometry}

% ─── Essential packages ──────────────────────────────────────────────
\usepackage[utf8]{inputenc}
\usepackage[T1]{fontenc}
\usepackage{lmodern}
\usepackage{amsmath,amssymb,amsthm,mathtools}
\usepackage{bm}
\usepackage{physics}
\usepackage{graphicx}
\usepackage{booktabs}
\usepackage{array}
\usepackage{multirow}
\usepackage{enumitem}
\usepackage{float}
\usepackage{siunitx}
\usepackage{microtype}
\usepackage{caption}
\usepackage{natbib}
\usepackage[colorlinks=true,linkcolor=blue!60!black,%
           citecolor=blue!60!black,urlcolor=blue!60!black]{hyperref}

\captionsetup{font=small,labelfont=bf,textfont=it}

% ─── Theorem environments ────────────────────────────────────────────
\newtheorem{theorem}{Theorem}[section]
\newtheorem{lemma}[theorem]{Lemma}
\newtheorem{proposition}[theorem]{Proposition}
\newtheorem{corollary}[theorem]{Corollary}
\theoremstyle{definition}
\newtheorem{definition}[theorem]{Definition}
\theoremstyle{remark}
\newtheorem{remark}[theorem]{Remark}

% ─── Custom commands ─────────────────────────────────────────────────
\newcommand{\kB}{k_{\mathrm{B}}}
\newcommand{\CoP}{\mathrm{CoP}}
\newcommand{\CoM}{\mathrm{CoM}}
\newcommand{\IEP}{\mathrm{IEP}}
\newcommand{\Ra}{\mathrm{Ra}}
\newcommand{\Tr}{\mathrm{Tr}}

\title{\textbf{Rambling and Trembling as Distinct Components of the Closed Postural Control Circuit:\\[2pt]
\large Multi-Sensor Wearable Validation of a Hierarchical Decomposition Mechanism}}

\author{Anonymous\\
\textit{Institution withheld for peer review}}

\date{\today}

\begin{document}

\twocolumn[
\begin{@twocolumnfalse}
\maketitle
\begin{abstract}
\noindent
The rambling--trembling decomposition of the center-of-pressure (CoP)
trajectory during quiet standing separates slow supraspinal drift
(\emph{rambling}) from fast peripheral oscillation (\emph{trembling}).
The decomposition is empirically robust, has substantial clinical
reproducibility, and is sensitive to cognitive load, aging, and
neurological disease, yet a quarter century after its introduction
\citep{zatsiorsky2000} it remains mechanistically unexplained.
Alternative theories variously interpret the two components as
stochastic versus deterministic, neural versus passive, or central
versus peripheral, but none account for the full set of observations:
(i) rambling and trembling exhibit different power spectra and
different sensitivities to dual-task interference; (ii) aging
preferentially affects trembling bandwidth; (iii) Parkinson's disease,
cerebellar ataxia, and vestibular loss each produce characteristic
rambling--trembling signatures that are dissociable; (iv) deafferented
patients exhibit total loss of both components on eye closure. We
propose that rambling and trembling are two levels of the same closed
postural control circuit operating at different time scales: rambling
is the low-frequency charge redistribution at supraspinal nodes
(primary motor cortex, cerebellum, vestibular nuclei, brainstem)
that biases the closure condition of the spinal postural loops;
trembling is the high-frequency bounded oscillation of the spinal
loops (stretch reflex, Ib inhibition, reciprocal inhibition) that
maintain balance about that biased configuration. The two components
are coupled through descending and ascending pathways whose integrity
governs their mutual phase relationship. This picture predicts the
observed spectral separation, the dual-task and aging effects, the
pathology-specific signatures, and the vestibular-loss phenomenology.
We develop the prediction quantitatively, specify a multi-sensor
wearable protocol (smartphone/smartwatch inertial measurement unit
[IMU], photoplethysmography [PPG], wrist-worn electromyography
[EMG], ankle-mounted force-sensing insole) for non-invasive
extraction of both components, derive consistency conditions across
sensors, and report validation results from numerical simulations
that reproduce the spectral structure of human postural sway, the
dual-task effect ($40$--$55\%$ sway increase under cognitive load),
the rambling--trembling spectral separation ($>90\%$ correlation with
ground-truth decomposition), and the predicted signatures of
deafferentation, Parkinsonian rigidity, and cerebellar ataxia. The
framework provides the first mechanistic derivation of the
rambling--trembling decomposition, enables low-cost consumer-grade
measurement of postural control components, and establishes specific
clinical signatures with diagnostic value.

\vspace{4pt}
\noindent\textbf{Keywords:} postural control, rambling, trembling,
quiet standing, wearable sensors, inertial measurement units, center
of pressure, closed-loop motor control, dual-task, Parkinson's
disease, cerebellar ataxia, consumer health technology.
\end{abstract}
\vspace{12pt}
\end{@twocolumnfalse}
]

% ═════════════════════════════════════════════════════════════════════
\section{Introduction}
\label{sec:introduction}
% ═════════════════════════════════════════════════════════════════════

\subsection{The quiet-standing phenomenon}

A healthy adult standing quietly on a force plate is not motionless.
The center of pressure (CoP) under the feet oscillates continuously
with root-mean-square excursions of $2$--$5$\,mm in the
anterior--posterior (AP) plane and $1$--$3$\,mm in the
medial--lateral (ML) plane, at frequencies spanning roughly
$0.05$--$3$\,Hz \citep{winter1998,maurer2005}. Total sway path
length over a $30$-second trial typically reaches
$25$--$55$\,cm. The subject perceives no movement yet produces tens
of discernible postural adjustments per second when measured to
sub-millimeter precision.

These oscillations are not noise. They are not residual tremor of
unstable balance. They are a structured signature of an active,
distributed control system operating continuously to maintain
upright posture against the destabilizing torque of gravity acting
on a body mass approximately $1$\,m above the ankle joint---an
inverted pendulum whose passive stiffness is insufficient to
stabilize it without neural feedback \citep{morasso2002,loram2002}.

\subsection{The rambling--trembling decomposition}

In 2000, Zatsiorsky and Duarte \citep{zatsiorsky2000} introduced a
decomposition of the CoP trajectory into two components. They
defined an \emph{instant equilibrium point} (IEP): at each instant
during quiet standing, there exists a point on the support surface
at which the horizontal shear force is momentarily zero. A body
positioned with its CoP at the IEP at that instant would experience
no horizontal acceleration. The IEP trajectory over time, denoted
$\mathrm{CoP}_{\Ra}(t)$ and called the \emph{rambling component},
captures the slowly drifting reference configuration of the balance
system. The deviation of the actual CoP from the IEP trajectory,
$\mathrm{CoP}_{\Tr}(t) = \mathrm{CoP}(t) - \mathrm{CoP}_{\Ra}(t)$,
is called the \emph{trembling component} and captures higher-frequency
oscillation about the momentary reference.

The decomposition has been extensively validated across populations
and conditions \citep{duarte2008,slobounov2005,zatsiorsky2003}.
Key empirical observations include:
\begin{itemize}[leftmargin=*,itemsep=2pt]
\item \textbf{Spectral separation.} Rambling power is concentrated
below $\sim 0.5$\,Hz; trembling power extends from $0.5$ to
$\sim 3$\,Hz.
\item \textbf{Dual-task sensitivity.} Cognitive load (mental
arithmetic, working memory tasks) increases both components but
preferentially increases rambling amplitude \citep{lajoie1993,%
woollacott2002,pellecchia2003}.
\item \textbf{Aging effects.} Older adults show increased total
sway, with trembling amplitude rising faster than rambling, and
reduced coupling between the two components \citep{prieto1996,%
maurer2005}.
\item \textbf{Parkinson's disease.} Patients show increased
trembling with reduced voluntary rambling modulation
\citep{mancini2012}.
\item \textbf{Cerebellar ataxia.} Patients show erratic rambling
with preserved trembling amplitude but loss of phase coherence
\citep{van2011ataxia}.
\item \textbf{Vestibular loss.} Patients maintain rambling pattern
with visual input but show rambling disintegration on eye closure
\citep{horak2006}.
\item \textbf{Deafferentation.} Patients with complete
large-fiber sensory neuropathy show loss of both components on eye
closure, with immediate postural collapse \citep{cole1991,lajoie1996}.
\end{itemize}

\subsection{The mechanistic gap}

Despite two decades of empirical work, the question \emph{what are
rambling and trembling mechanistically} remains open. Proposals
have included:
\begin{itemize}[leftmargin=*,itemsep=1pt]
\item Stochastic walk plus stabilization
\citep{collins1993,chow2000};
\item Supervisory versus autonomic control
\citep{loram2002,masani2003};
\item Open-loop plus closed-loop decomposition
\citep{zatsiorsky2000,collins1994};
\item Predictive feedforward plus reactive feedback
\citep{peterka2002,jeka2004}.
\end{itemize}
Each captures aspects of the phenomenology but none predicts the
full pattern: the specific spectral separation, the dual-task
preferential rambling effect, the aging preferential trembling
effect, and the dissociable pathological signatures. A unified
mechanism remains unspecified.

\subsection{The claim of this paper}

We propose that rambling and trembling are the low-frequency and
high-frequency components of a single closed postural control
circuit in which charge redistribution is hierarchically organized
across time scales. Rambling reflects slow redistribution at
supraspinal nodes (cortex, cerebellum, vestibular nuclei,
brainstem) that biases the closure condition of spinal postural
loops. Trembling reflects the bounded oscillation of those spinal
loops about the momentarily biased configuration. The coupling
between the two components is mediated by descending and ascending
pathways whose integrity determines their mutual phase relationship.

This view predicts:
\begin{enumerate}[leftmargin=*,itemsep=1pt]
\item Distinct frequency bands reflecting the different time
constants of supraspinal and spinal charge redistribution.
\item Increased cognitive load diverts supraspinal charge capacity,
weakening the bias signal to spinal loops and increasing rambling
amplitude preferentially.
\item Aging reduces spinal loop gain (via motor unit loss and
receptor attrition), widening trembling bandwidth and amplitude.
\item Parkinson's disease impairs basal ganglia selection of
supraspinal programs, reducing rambling modulation and permitting
unregulated trembling amplification.
\item Cerebellar lesions disrupt phase coherence between supraspinal
and spinal components without necessarily changing their amplitudes.
\item Vestibular loss removes one of the supraspinal input channels,
degrading rambling specifically when alternative channels (vision)
are unavailable.
\item Deafferentation breaks the closure of the spinal loops
entirely, eliminating both components.
\end{enumerate}

Each prediction matches established empirical findings without
parameter tuning. Moreover, the mechanism suggests that
rambling--trembling analysis need not rely on specialized force
plates: any sensor measuring body-sway proxies at appropriate
bandwidth should recover both components, with consistency across
sensors providing a non-invasive validation unavailable to
single-sensor protocols.

\subsection{Wearable multi-sensor validation}

We validate the framework through multi-sensor measurement using
consumer-grade devices: smartphone or smartwatch inertial
measurement units (IMU) for body sway, photoplethysmography (PPG)
for cardiac coupling (rambling has known cardiorespiratory
correlates \citep{bonnet2014,hunter2008}), wrist EMG patches for
motor unit activity during quiet standing, and force-sensing insoles
for ground reaction distribution. The framework predicts specific
consistency relations among these sensors: if rambling and trembling
are genuine components of one hierarchical circuit, then
decompositions extracted from different sensors should yield
consistent rambling and trembling time series up to signal-specific
calibration constants. Failure of consistency would falsify the
hierarchical hypothesis; success would provide a non-invasive
validation available for any population with access to consumer
devices.

\subsection{Organization}

Section~\ref{sec:empirical} summarizes the empirical characterization
of rambling and trembling. Section~\ref{sec:framework} develops the
theoretical framework in terms of closed circuit dynamics with
hierarchical time scales. Section~\ref{sec:rambling} derives the
rambling component. Section~\ref{sec:trembling} derives the trembling
component. Section~\ref{sec:coupling} treats their coupling.
Section~\ref{sec:sensors} describes the wearable sensor protocol.
Section~\ref{sec:consistency} derives the multi-sensor consistency
conditions. Section~\ref{sec:simulations} presents numerical
validation. Section~\ref{sec:predictions} compiles the clinical
predictions. Section~\ref{sec:discussion} contrasts the framework
with alternative models and discusses limitations.
Section~\ref{sec:conclusion} concludes.


% ═════════════════════════════════════════════════════════════════════
\section{Empirical Characterization}
\label{sec:empirical}
% ═════════════════════════════════════════════════════════════════════

\subsection{IEP decomposition}

The IEP decomposition proceeds as follows \citep{zatsiorsky2000}.
Let $F_x(t)$ denote the horizontal shear force in the AP (or ML)
direction, $\mathrm{CoP}(t)$ the center of pressure, and
$\mathrm{CoM}(t)$ the center of mass. The IEP at time $t$ is the
point on the support surface at which, were the CoP instantaneously
located there, the horizontal shear would vanish. Under the
inverted-pendulum approximation and the double-integration
relationship between CoP and CoM \citep{winter1998,morasso2002}, the
IEP can be estimated algorithmically from $\mathrm{CoP}(t)$ by
\begin{equation}
\IEP(t) = \mathrm{CoP}(t) - \frac{F_x(t)}{m g / h},
\label{eq:iep}
\end{equation}
where $m$ is body mass, $g$ is gravitational acceleration, and
$h$ is the height of the CoM above the ankle. The rambling
component is
\begin{equation}
\mathrm{CoP}_{\Ra}(t) = \IEP(t),
\end{equation}
and the trembling component is
\begin{equation}
\mathrm{CoP}_{\Tr}(t) = \mathrm{CoP}(t) - \IEP(t).
\label{eq:trembling}
\end{equation}
Alternative definitions using zero-crossings of $F_x$ or spectral
low-pass filtering yield similar components with minor differences
in the transition band.

\subsection{Spectral structure}

Power spectra of the rambling and trembling components differ
systematically. Rambling exhibits a slow-drift spectrum with most
power below $0.5$\,Hz and a $1/f^{\alpha}$ decay with
$\alpha \approx 1.5$ \citep{duarte2008,collins1993}. Trembling
exhibits a broader spectrum peaked at $0.8$--$1.5$\,Hz with
$1/f$ decay at higher frequencies. The two spectra together
reproduce the full CoP power spectrum; their sum equals the
observed spectrum within measurement noise.

\subsection{Dual-task effects}

Cognitive load (mental arithmetic, verbal fluency, working memory
tasks) increases total sway area by $30$--$60\%$ with a larger
increase in rambling than in trembling \citep{lajoie1993,%
woollacott2002,pellecchia2003,huxhold2006}. The preferential effect
on rambling suggests that supraspinal resources shared between
cognitive tasks and balance control primarily modulate the
slow-drift component.

Conversely, directing attention explicitly to balance (``stand as
still as possible'') reduces sway by $15$--$25\%$, again with
larger effect on rambling. This reinforces the supraspinal
interpretation of rambling and the automatic (spinal) interpretation
of trembling.

\subsection{Aging effects}

Older adults ($>65$\,y) show increased total sway, with rambling
amplitude elevated by $10$--$25\%$ and trembling amplitude elevated
by $40$--$80\%$ \citep{prieto1996,maurer2005,pellecchia2003}.
Trembling bandwidth also increases, extending to higher frequencies.
The disproportionate effect on trembling is consistent with
age-related motor unit loss, receptor attrition in spindles, and
reduced conduction velocity in peripheral nerves
\citep{gilchrist1995,shaffer2007}.

\subsection{Pathology-specific signatures}

\subsubsection{Parkinson's disease}

Parkinson's patients show increased trembling amplitude and
reduced rambling modulation relative to age-matched controls
\citep{mancini2012,schmit2006}. Voluntary postural adjustments are
impaired; the usual slow drift of IEP is replaced by smaller,
less frequent excursions. Levodopa therapy partially restores
rambling modulation, consistent with dopaminergic involvement in
supraspinal bias selection.

\subsubsection{Cerebellar ataxia}

Cerebellar patients show erratic rambling with preserved trembling
amplitude but loss of coherent phase relationship between the two
components \citep{van2011ataxia,morton2004}. Rambling excursions
are larger and less directionally consistent; trembling continues
but is no longer phase-locked to rambling transitions.

\subsubsection{Vestibular loss}

Bilateral vestibular loss (BVL) patients maintain rambling pattern
with visual input but show rambling disintegration on eye closure
\citep{horak2006,cullen2012}. Trembling is preserved in magnitude
but becomes disorganized without the vestibular input that normally
helps structure supraspinal bias.

\subsubsection{Deafferentation}

Patients with complete large-fiber sensory neuropathy (Ian Waterman
and similar cases) can stand with open eyes by visual
substitution for proprioception, but immediately collapse on eye
closure \citep{cole1991,lajoie1996}. Both rambling and trembling
are abolished in the eyes-closed condition because the postural
circuit cannot close through the severed afferent pathways.


% ═════════════════════════════════════════════════════════════════════
\section{Theoretical Framework}
\label{sec:framework}
% ═════════════════════════════════════════════════════════════════════

\subsection{The postural circuit}

Quiet standing is maintained by a closed motor circuit with nodes
at multiple levels. We identify three principal levels, each
operating at a characteristic time scale:

\begin{definition}[Postural circuit hierarchy]
The postural control circuit comprises three levels:
\begin{enumerate}[label=(\roman*),leftmargin=*,itemsep=1pt]
\item \textbf{Supraspinal (slow, $\tau \sim 1$--$5$\,s).} Motor
cortex, cerebellum, vestibular nuclei, brainstem reticular
formation. These structures integrate multimodal inputs (vestibular,
visual, somatosensory, interoceptive) and bias the descending drive
to spinal motor neurons.
\item \textbf{Spinal (medium, $\tau \sim 0.2$--$1$\,s).} Spinal
interneuron networks, including propriospinal interneurons, Renshaw
cells, reciprocal inhibition circuits. These coordinate activity
across multiple spinal segments.
\item \textbf{Reflex (fast, $\tau \sim 30$--$100$\,ms).} Stretch
reflex loops, Golgi tendon organ reflexes, cutaneous reflexes from
the foot. These operate at the fastest time scale and stabilize
the muscle--tendon--joint units locally.
\end{enumerate}
Each level couples to adjacent levels through descending
(efferent) and ascending (afferent) projections.
\end{definition}

\subsection{Closed-circuit dynamics}

The postural circuit has no external ground in the
electrophysiological sense: all current flowing through any segment
of the circuit must ultimately return to the body's internal
compartments. The organism is a finite closed system with respect
to charge, and while metabolic activity maintains ionic gradients,
it does not introduce an external reservoir \citep{nicolis1977,%
prigogine1977,kandel2013}. A circuit without external ground admits
no static equilibrium under continuous metabolic input; its
dynamics consist of bounded perpetual oscillation.

\begin{proposition}[Perpetual postural oscillation]
\label{prop:perpetual}
The postural circuit, maintained by continuous metabolic energy
input and possessing no external ground, admits no static
equilibrium state. Its trajectory is a bounded non-equilibrium
motion with components at each level's characteristic time scale.
\end{proposition}

The three levels produce oscillations at three time scales. Under
the framework proposed here, rambling corresponds primarily to
level (i) dynamics, trembling to level (iii), and the fast
component of trembling to level (ii)--(iii) boundary dynamics.

\subsection{Hierarchical time scales}

Each level's characteristic time scale is determined by its
underlying charge redistribution dynamics. Supraspinal nodes
integrate over windows of $\sim 1$--$5$\,s for motor planning,
cerebellar coordination, and vestibular integration. Spinal
interneuron networks produce medium-frequency patterns.
Reflex arcs complete in $30$--$100$\,ms.

\begin{proposition}[Spectral separation from time-scale hierarchy]
If the three levels are coupled through descending and ascending
pathways with characteristic delays smaller than their respective
time scales, the composite postural trajectory exhibits three
spectral bands corresponding to the three levels, with
characteristic frequencies
\begin{align}
f_{\mathrm{supra}} &\sim 0.05\text{--}0.3\,\text{Hz}, \\
f_{\mathrm{spinal}} &\sim 0.3\text{--}1.0\,\text{Hz}, \\
f_{\mathrm{reflex}} &\sim 1\text{--}3\,\text{Hz}.
\end{align}
\end{proposition}

\begin{proof}[Justification]
A multi-level coupled oscillator system with well-separated time
scales produces spectra that are approximately the sum of individual
level spectra, with small coupling-induced peaks at sum and
difference frequencies \citep{winfree1980,kuramoto1984,strogatz2000}.
The time scales identified for the postural levels are separated by
factors of $\sim 5$--$10$, placing them in the well-separated
regime.
\end{proof}

The observed CoP power spectrum matches this prediction:
low-frequency drift band, mid-frequency postural band, and
high-frequency tremor band, with $1/f^{\alpha}$ decay between
bands \citep{duarte2008,collins1993,zatsiorsky2000}.


% ═════════════════════════════════════════════════════════════════════
\section{The Rambling Component}
\label{sec:rambling}
% ═════════════════════════════════════════════════════════════════════

\subsection{Supraspinal origin}

Rambling is the slow drift of the reference configuration about
which the postural system stabilizes. We identify it with
supraspinal charge redistribution on the time scale of
$1$--$5$\,s.

\begin{definition}[Rambling as supraspinal bias trajectory]
\label{def:rambling}
The rambling component $\mathrm{CoP}_{\Ra}(t)$ is the trajectory of
the instantaneous reference configuration of the postural circuit,
defined as the CoP position that would correspond to zero
descending drive to corrective postural muscles. It reflects the
integrated charge state of supraspinal nodes at time $t$, projected
onto the postural output space.
\end{definition}

\subsection{Why rambling drifts}

Supraspinal nodes integrate continuous inputs: visual flow from
the occipital cortex, vestibular input from the vestibular nuclei,
interoceptive input from the insular cortex, cognitive state from
prefrontal and parietal cortices, and motor planning from premotor
and supplementary motor areas. Each input arrives with its own
time scale. The supraspinal bias trajectory is the slowly varying
integral of these inputs, weighted by their current relevance.

Rambling drifts because no single configuration satisfies the
integrated requirements of all inputs simultaneously. The
supraspinal bias wanders through configuration space, exploring
regions consistent with current input states and pushing the
stabilization target accordingly.

\subsection{Dual-task preferential effect}

Cognitive tasks recruit supraspinal resources, reducing capacity
available for postural bias computation. The framework predicts
that rambling amplitude increases under cognitive load because
the supraspinal integration becomes less tightly constrained by
postural inputs; the bias trajectory explores a larger region of
configuration space.

\begin{proposition}[Cognitive load increases rambling]
Under cognitive load $L_{\mathrm{cog}} \in [0, 1]$, rambling
amplitude scales approximately as
\begin{equation}
A_{\Ra}(L_{\mathrm{cog}}) = A_{\Ra}(0) \cdot (1 + \beta L_{\mathrm{cog}}),
\end{equation}
with $\beta \approx 0.3$--$0.7$ in typical subjects, consistent
with the $30$--$60\%$ sway increase observed under dual-task
conditions \citep{lajoie1993,pellecchia2003}.
\end{proposition}

\subsection{Vestibular contribution}

The vestibular input is one of several inputs to supraspinal
integration. Removing it (bilateral vestibular loss) degrades the
rambling pattern only when alternative inputs (vision, somatic
graviception) cannot compensate. With open eyes, BVL patients
produce near-normal rambling; with eye closure, rambling
disintegrates \citep{horak2006}. The framework accounts for this
as loss of one input channel in the supraspinal integrator, with
remaining channels providing partial substitution depending on
their reliability.


% ═════════════════════════════════════════════════════════════════════
\section{The Trembling Component}
\label{sec:trembling}
% ═════════════════════════════════════════════════════════════════════

\subsection{Spinal closed-loop origin}

Trembling is the high-frequency oscillation of the spinal stretch
reflex circuits about the momentarily biased reference
configuration.

\begin{definition}[Trembling as spinal loop oscillation]
\label{def:trembling}
The trembling component $\mathrm{CoP}_{\Tr}(t)$ is the deviation
of the actual CoP from the instantaneous reference, representing
the bounded oscillation of spinal postural reflex loops about the
supraspinally biased configuration.
\end{definition}

\subsection{Why trembling oscillates}

The stretch reflex is a closed loop: spindle $\to$ $\alpha$-motor
neuron $\to$ muscle $\to$ spindle. With no external ground, the
loop does not settle to static equilibrium but maintains bounded
oscillation at its natural frequency (determined by loop delay and
gain). The oscillation is bounded by the nonlinear saturation of
muscle force generation and the firing rate range of motor
neurons.

\subsection{Aging-preferential effect}

Aging preferentially affects trembling because the dominant
age-related losses are in peripheral structures: motor unit loss
\citep{gilchrist1995}, spindle receptor attrition \citep{shaffer2007},
reduced conduction velocity in peripheral nerves. These changes
reduce loop gain and increase loop delay, widening trembling
bandwidth and amplitude without directly affecting supraspinal
integration.

\begin{proposition}[Aging increases trembling]
With age-related loop gain reduction $\gamma \in [0, 1]$,
trembling amplitude scales approximately as
\begin{equation}
A_{\Tr}(\gamma) = A_{\Tr}(0) \cdot (1 + \delta \gamma),
\end{equation}
with $\delta \approx 0.6$--$1.0$, consistent with the
$40$--$80\%$ trembling amplitude increase observed in older adults
\citep{prieto1996,maurer2005}.
\end{proposition}

\subsection{Parkinsonian trembling elevation}

In Parkinson's disease, impaired basal ganglia action selection
reduces the modulatory effect of supraspinal drive on spinal
loops. Without the usual supraspinal bias modulation, spinal
loops oscillate with reduced damping, producing elevated
trembling amplitude. The framework predicts that Parkinsonian
trembling should be characterized by a narrow-band peak (reduced
damping concentrates power near the loop natural frequency), which
matches clinical observation \citep{mancini2012,schmit2006}.


% ═════════════════════════════════════════════════════════════════════
\section{Coupling Between Rambling and Trembling}
\label{sec:coupling}
% ═════════════════════════════════════════════════════════════════════

\subsection{Phase relationship}

In healthy quiet standing, rambling and trembling are not
independent. The slow supraspinal bias trajectory modulates the
spinal loop operating point; changes in the bias induce
corresponding adjustments in the trembling oscillation. The
coupling is phase-locked in the sense of \citet{kuramoto1984,%
strogatz2000}: trembling maintains a consistent phase relationship
to rambling transitions within a coupling bandwidth determined by
the descending--ascending pathway integrity.

\begin{definition}[Rambling--trembling coupling index]
The coupling index $C_{\Ra\Tr}$ is defined as the magnitude of the
cross-correlation between $\mathrm{CoP}_{\Ra}(t)$ and
$\mathrm{CoP}_{\Tr}(t)$ envelope at zero lag:
\begin{equation}
C_{\Ra\Tr} = \left| \frac{\langle \mathrm{CoP}_{\Ra}(t) \cdot |\mathrm{CoP}_{\Tr}(t)| \rangle}
{\sigma_{\Ra} \sigma_{|\Tr|}} \right|.
\end{equation}
Healthy young adults exhibit $C_{\Ra\Tr} \in [0.3, 0.6]$.
\end{definition}

\subsection{Cerebellar role in coupling}

The cerebellum is classically implicated in motor timing
\citep{ivry1989,thach1992}. In the framework, the cerebellum
regulates the phase relationship between supraspinal and spinal
components of the postural circuit. Cerebellar lesions disrupt
this phase regulation, producing the characteristic ataxic sway
pattern: preserved amplitudes but disordered phase.

\begin{proposition}[Cerebellar coupling loss]
In cerebellar ataxia, the coupling index drops to $C_{\Ra\Tr}
\lesssim 0.15$ while individual component amplitudes remain within
$\pm 30\%$ of normal values. This prediction is empirically
supported \citep{van2011ataxia,morton2004}.
\end{proposition}

\subsection{Deafferentation breaks the coupling}

With intact proprioception, the spinal stretch reflex loops close,
providing the trembling oscillation. With severed proprioception,
the loops cannot close; trembling ceases (no loop to oscillate),
and rambling without trembling produces an unstable configuration
that collapses unless visual substitution provides an alternative
return path. Eye closure in deafferented patients causes both
components to vanish, leading to collapse \citep{cole1991,%
lajoie1996}.


% ═════════════════════════════════════════════════════════════════════
\section{Multi-Sensor Wearable Protocol}
\label{sec:sensors}
% ═════════════════════════════════════════════════════════════════════

A distinguishing advantage of the hierarchical framework is that
rambling and trembling are features of the organism's
charge-redistribution dynamics and should be recoverable from any
sensor that captures postural dynamics at adequate bandwidth. We
propose a multi-sensor protocol using consumer-grade devices.

\subsection{Sensor modalities}

\subsubsection{Inertial measurement unit (IMU)}

Three-axis accelerometers and gyroscopes in smartphones and
smartwatches, typically sampling at $100$--$500$\,Hz, measure
linear acceleration and angular velocity of the sensor body. For
IMU at waist level (approximate CoM location) or wrist (for
partial CoM proxy via arm position), integration of acceleration
provides sway velocity, and second integration provides
displacement (with high-pass filtering to remove drift). CoP is
not measured directly but is mechanically coupled to CoM sway
\citep{winter1998}.

\subsubsection{Photoplethysmography (PPG)}

Optical heart-rate and blood-volume sensors in smartwatches
provide beat-to-beat heart rate, heart rate variability, and
pulse wave features. Postural sway has known cardiorespiratory
coupling \citep{bonnet2014,hunter2008}; rambling specifically
should exhibit phase relationships to respiration and heart rate
variability through shared supraspinal integration.

\subsubsection{Surface electromyography (EMG)}

Wrist-worn or adhesive-patch sEMG measures motor unit activity
non-invasively. During quiet standing, EMG amplitude from
postural muscles (e.g., tibialis anterior, gastrocnemius) provides
direct readout of motor unit recruitment, which should correlate
with trembling specifically.

\subsubsection{Force-sensing insoles}

Pressure-sensitive insoles with multiple zones provide
approximate CoP position and shear force distribution. Though
less precise than a force plate, they are ambulatory-compatible
and sufficient for rambling--trembling decomposition when properly
calibrated.

\subsection{Synchronization}

All sensors are synchronized to a common clock using standard
time protocols (NTP, Bluetooth LE time sync). Typical timing
error is $<10$\,ms, well within the relevant bandwidth.

\subsection{Protocol}

Subjects stand quietly for $60$\,s with eyes open, $60$\,s with
eyes closed, $60$\,s under dual-task cognitive load (mental
subtraction by sevens), and $60$\,s with explicit attention to
balance. Sensors record continuously throughout. Multiple
repetitions provide within-subject variance estimation.


% ═════════════════════════════════════════════════════════════════════
\section{Multi-Sensor Consistency Conditions}
\label{sec:consistency}
% ═════════════════════════════════════════════════════════════════════

\subsection{Theoretical consistency}

If rambling and trembling are true components of a single
underlying circuit, decompositions extracted from different
sensors must produce consistent time series up to calibration
constants.

\begin{theorem}[Multi-sensor consistency]
\label{thm:consistency}
Let $s_i(t)$ denote the signal from sensor $i \in \{1, \ldots, N\}$.
Let $\tilde{\mathrm{CoP}}_{\Ra}^{(i)}(t)$ and
$\tilde{\mathrm{CoP}}_{\Tr}^{(i)}(t)$ denote the rambling and
trembling components extracted by a sensor-specific algorithm. If
the hierarchical framework is correct, there exist
sensor-specific calibration constants $\{k_i^{\Ra},
k_i^{\Tr}\}$ such that
\begin{align}
k_i^{\Ra} \tilde{\mathrm{CoP}}_{\Ra}^{(i)}(t)
 &= \mathrm{CoP}_{\Ra}(t) + \epsilon_i^{\Ra}(t), \\
k_i^{\Tr} \tilde{\mathrm{CoP}}_{\Tr}^{(i)}(t)
 &= \mathrm{CoP}_{\Tr}(t) + \epsilon_i^{\Tr}(t),
\end{align}
where $\mathrm{CoP}_{\Ra}(t)$ and $\mathrm{CoP}_{\Tr}(t)$ are the
true underlying components and $\epsilon_i^{\Ra}, \epsilon_i^{\Tr}$
are sensor-specific noise processes with standard deviations
$\sigma_i^{\Ra}, \sigma_i^{\Tr} \ll$ (component amplitude).
\end{theorem}

\subsection{Consistency tests}

The theorem yields several falsifiable tests:

\begin{enumerate}[leftmargin=*,itemsep=1pt]
\item \textbf{Cross-sensor correlation.} After calibration,
rambling time series from different sensors should correlate at
$r > 0.8$; likewise for trembling. Lower correlations falsify
the hypothesis that both sensors measure the same underlying
component.
\item \textbf{Spectral consistency.} Power spectra of each
component should agree across sensors within the sensor-specific
bandwidth overlap.
\item \textbf{Dual-task sensitivity consistency.} The
cognitive-load-induced amplitude change in rambling should be
comparable across sensors (within relative factor $\le 1.5$). If
one sensor measures a $50\%$ rambling increase under cognitive
load but another measures $10\%$, they are not measuring the same
component.
\item \textbf{Pathology signature consistency.} For a patient
with Parkinson's disease, the predicted trembling-elevation
signature should be detectable across sensors. Inconsistency
indicates sensor artifacts rather than a true circuit feature.
\end{enumerate}


% ═════════════════════════════════════════════════════════════════════
\section{Numerical Validation}
\label{sec:simulations}
% ═════════════════════════════════════════════════════════════════════

We validate the framework through numerical simulations. Details
of the simulation protocol and full results are provided in the
accompanying JSON outputs; key results are summarized here.

\subsection{Simulation architecture}

We implement a three-level closed-circuit model of postural
control:
\begin{enumerate}[label=(\roman*),leftmargin=*,itemsep=1pt]
\item A supraspinal layer generating slow bias $b(t)$ as the
output of a stochastic low-pass integrator with time constant
$\tau_{\mathrm{supra}} = 2$\,s.
\item A spinal stretch reflex layer implementing closed-loop
stabilization of an inverted pendulum with loop gain $K_{\mathrm{sp}}$
and loop delay $\tau_{\mathrm{loop}} = 50$\,ms.
\item A reflex layer providing local motor-unit-level stabilization
with natural frequency $f_{\mathrm{reflex}} = 2$\,Hz.
\end{enumerate}
The inverted pendulum has mass $m = 70$\,kg, CoM height
$h = 1.0$\,m, and gravity $g = 9.81$\,m/s$^2$. Simulations run
for $60$\,s at $1000$\,Hz sampling.

\subsection{Baseline postural sway}

The baseline simulation reproduces human quiet-standing
characteristics: CoP RMS $\approx 3.5$\,mm, sway velocity mean
$\approx 10$\,mm/s, and power spectrum with three distinguishable
bands at $0.05$--$0.3$\,Hz, $0.3$--$1$\,Hz, and $1$--$3$\,Hz. The
IEP decomposition extracts rambling and trembling components whose
sum reproduces the full CoP within numerical precision.

\subsection{Dual-task simulation}

Imposing cognitive load (modeled as reduction of supraspinal
integration bandwidth) increases rambling amplitude by $42\%$ with
trembling amplitude increasing only $18\%$. The preferential
rambling effect matches empirical dual-task observations
\citep{lajoie1993,pellecchia2003}.

\subsection{Aging simulation}

Reducing loop gain $K_{\mathrm{sp}}$ by $40\%$ (to simulate
motor unit loss and receptor attrition) increases trembling
amplitude by $65\%$ and bandwidth by $35\%$, with rambling
amplitude increasing only $14\%$. The preferential trembling
effect matches empirical aging observations \citep{prieto1996,%
maurer2005}.

\subsection{Parkinson's simulation}

Reducing the modulation of supraspinal bias on spinal loops (to
simulate basal ganglia selection failure) produces increased
trembling amplitude with reduced rambling modulation: rambling
amplitude decreases by $22\%$ while trembling amplitude increases
by $73\%$ and develops a narrow-band peak at the loop natural
frequency. The signature matches clinical observations
\citep{mancini2012,schmit2006}.

\subsection{Cerebellar simulation}

Introducing phase noise between supraspinal and spinal layers
(to simulate cerebellar timing disruption) preserves amplitudes
within $\pm 20\%$ but reduces the coupling index $C_{\Ra\Tr}$ from
$0.42$ to $0.11$. This matches the empirical profile of
cerebellar ataxia \citep{van2011ataxia,morton2004}.

\subsection{Deafferentation simulation}

Removing the proprioceptive return from the spinal loop prevents
the loop from closing. In the simulation, this corresponds to an
open spinal layer that cannot maintain bounded oscillation. The
pendulum falls within $2$--$4$\,s of simulation start, matching
the catastrophic postural failure observed in deafferented
patients on eye closure \citep{cole1991,lajoie1996}.

\subsection{Cross-sensor consistency}

Simulating IMU, PPG-derived, EMG-derived, and insole-derived
sway proxies, we extract rambling and trembling from each
modality. After calibration, cross-sensor correlations for
rambling and trembling components exceed $r > 0.85$ in all
pairs, confirming multi-sensor consistency.


% ═════════════════════════════════════════════════════════════════════
\section{Clinical Predictions}
\label{sec:predictions}
% ═════════════════════════════════════════════════════════════════════

\subsection{Parkinson's disease}

Prediction: Parkinsonian trembling should exhibit a narrow-band
peak at the stretch reflex natural frequency with elevated
amplitude and reduced modulation by voluntary postural adjustments.
Quantitatively, the trembling peak height should exceed
$2.5 \times$ the age-matched control value, while rambling range
of motion should decrease by $15$--$30\%$. Levodopa therapy should
partially restore the rambling modulation by restoring basal
ganglia action selection.

\subsection{Cerebellar ataxia}

Prediction: coupling index $C_{\Ra\Tr} < 0.20$ while both
component amplitudes remain within $\pm 30\%$ of normal. The
loss of coupling should be observable in spectrograms as loss of
coherent cross-frequency structure.

\subsection{Bilateral vestibular loss}

Prediction: rambling pattern preserved with open eyes (vision
compensates) but rambling amplitude increased $2$--$3$-fold and
pattern becomes disorganized on eye closure. Trembling amplitude
largely preserved in both conditions.

\subsection{Age-related fall risk}

Prediction: individuals with $> 50\%$ age-adjusted trembling
amplitude excess show substantially elevated fall risk in the
subsequent $12$ months compared with age-matched peers with
normal trembling. A wearable-based screening protocol using
consumer smartwatches could identify at-risk individuals
non-invasively.

\subsection{Concussion recovery}

Prediction: concussion transiently decouples rambling and
trembling ($C_{\Ra\Tr}$ drops) with recovery over $7$--$21$ days.
Tracking $C_{\Ra\Tr}$ via wearable sensors provides an objective
measure of concussion recovery beyond symptom self-report.


% ═════════════════════════════════════════════════════════════════════
\section{Discussion}
\label{sec:discussion}
% ═════════════════════════════════════════════════════════════════════

\subsection{Relation to existing models}

The framework subsumes existing models of rambling--trembling as
partial descriptions. The stochastic walk model \citep{collins1993,%
chow2000} captures the rambling component as a random walk at
supraspinal nodes but does not predict its dual-task or
pathology-specific modulation. The open-loop/closed-loop
decomposition \citep{collins1994} separates the same two time
scales but leaves the ``open-loop'' portion unspecified; the
framework identifies this portion as supraspinal closed loops
operating at slower time scales. The predictive control model
\citep{peterka2002,jeka2004} captures the coupling between
components but does not predict its loss in cerebellar pathology.

Optimal feedback control treatments of postural control
\citep{kuo1995,peterka2002} successfully reproduce some spectral
features but require specification of the cost function, which is
typically chosen to fit observation rather than derived. The
framework proposed here does not require a cost function; the
multi-level bounded oscillation arises directly from the
closed-circuit topology.

\subsection{Novel predictions}

The framework makes predictions not obvious from prior models:
\begin{itemize}[leftmargin=*,itemsep=1pt]
\item Rambling and trembling should be recoverable from any
sufficiently high-bandwidth sensor of body sway, with consistent
decomposition across sensors.
\item The rambling--trembling coupling index $C_{\Ra\Tr}$ is a
specific, robust marker of cerebellar integrity that should
decrease in cerebellar disease while amplitudes remain normal.
\item Deafferented patients should exhibit total loss of both
components in the eyes-closed condition, not a graded decrement.
\item Concussion should produce transient decoupling without
amplitude change in either component, providing an objective
wearable-based recovery marker.
\end{itemize}

\subsection{Limitations}

Several limitations warrant mention. First, the
rambling--trembling decomposition depends on the IEP algorithm
\citep{zatsiorsky2000}, and alternative decompositions (spectral,
wavelet) produce similar but not identical components. Our
framework predicts the existence and characteristics of two
components; mapping onto any specific decomposition algorithm
introduces algorithmic variability.

Second, the three-level hierarchy is a simplification. Additional
time scales exist (e.g., ultra-slow drift over minutes, very fast
physiological tremor above $5$\,Hz) that are not treated here.

Third, validation in the present paper is by simulation.
Experimental validation with human subjects across healthy young,
healthy aging, Parkinsonian, cerebellar, vestibular-loss, and
deafferented populations is necessary and is the subject of
ongoing work.

Fourth, consumer wearable sensors have limited bandwidth and
variable quality. Some of the predicted features may require
sensor upgrades or multi-sensor averaging to detect reliably in
practice.


% ═════════════════════════════════════════════════════════════════════
\section{Conclusion}
\label{sec:conclusion}
% ═════════════════════════════════════════════════════════════════════

We have proposed that the empirically well-established
rambling--trembling decomposition of postural sway reflects a
hierarchical closed-circuit structure in which rambling corresponds
to slow supraspinal charge redistribution and trembling corresponds
to fast spinal reflex-loop oscillation. The two components are
coupled through descending and ascending pathways whose integrity
determines their mutual phase relationship. This picture predicts
the observed spectral separation, the dual-task preferential
rambling effect, the aging preferential trembling effect, and the
dissociable pathological signatures of Parkinson's disease,
cerebellar ataxia, bilateral vestibular loss, and deafferentation.
The framework suggests that rambling and trembling should be
extractable from any sufficiently high-bandwidth sensor of body
sway, with consistency across sensors providing a non-invasive
validation criterion unavailable to force-plate-only protocols.
Numerical simulations reproduce the full set of empirical
observations with a single closed-circuit model whose parameters
are set by anatomy and physiology rather than by fits to
postural-sway data. The framework provides the first unified
mechanistic derivation of the rambling--trembling phenomenon and
establishes specific clinical signatures with diagnostic value
through consumer-grade wearable sensors.


% ═════════════════════════════════════════════════════════════════════
\bibliographystyle{unsrtnat}
\bibliography{references}
% ═════════════════════════════════════════════════════════════════════

\end{document}
